\documentclass[conference]{IEEEtran}
\IEEEoverridecommandlockouts

\usepackage{cite}
\usepackage{amsmath,amssymb,amsfonts,amsthm}
\usepackage{algorithmic}
\usepackage{graphicx}
\usepackage{textcomp}
\usepackage{xcolor}
\usepackage{booktabs}
\usepackage{bm}
\usepackage{array}
\usepackage{url}

\newtheorem{definition}{Definition}
\newtheorem{theorem}{Theorem}
\newtheorem{proposition}{Proposition}
\newtheorem{remark}{Remark}

\def\BibTeX{{\rm B\kern-.05em{\sc i\kern-.025em b}\kern-.08em
    T\kern-.1667em\lower.7ex\hbox{E}\kern-.125emX}}

\begin{document}

\title{Gaussian Process Models for Spatial Prediction:\\
A Comprehensive Survey}

\author{\IEEEauthorblockN{Surya Rao Rayarao}
\IEEEauthorblockA{\textit{Department of Statistics and Data Sciences} \\
\textit{Department of Computer Science} \\
\textit{The University of Texas at Austin}\\
Austin, TX, USA \\
suryarao.r@utexas.edu}
\and
\IEEEauthorblockN{Naga Donikena}
\IEEEauthorblockA{\textit{Department of Statistics and Data Sciences} \\
\textit{Department of Computer Science} \\
\textit{The University of Texas at Austin}\\
Austin, TX, USA \\
naga.donikena@utexas.edu}
}

\maketitle

\begin{abstract}
Gaussian process (GP) models provide a principled probabilistic framework for spatial prediction, enabling inference over continuously indexed spatial fields while naturally encoding dependence structure through covariance functions. This survey reviews the mathematical foundations, key algorithms, and practical applications of GP-based spatial modeling, with emphasis on their role in data science tasks involving structured, dependent observations. We cover the theoretical underpinnings of spatial stationarity, variogram modeling, and kriging, alongside modern computational approaches including sparse approximations, inducing-point methods, and hierarchical GP frameworks. The paper systematically addresses four core challenges: identifying spatial dependencies in structured data, applying appropriate models that respect those dependencies, producing calibrated predictions that account for spatial structure, and recovering latent spatial fields from noisy or indirect observations. We survey representative applications in environmental monitoring, geostatistics, climate science, and machine learning, and discuss connections to related methodologies such as Gaussian Markov random fields, deep GPs, and neural process models. Computational limitations, stationarity assumptions, and scalability challenges are critically examined. The survey is intended for practitioners in data science and statistics seeking a rigorous yet accessible treatment of GP methods for spatial data.
\end{abstract}

\begin{IEEEkeywords}
Gaussian processes, spatial statistics, kriging, covariance functions, geostatistics, sparse approximations, spatial prediction, latent variable models
\end{IEEEkeywords}

%=============================================================================
\section{Introduction}
%=============================================================================

Spatial data arise throughout science, engineering, and industry whenever measurements are collected at geographically or geometrically referenced locations. Environmental sensor networks record temperature, precipitation, and pollutant concentrations across continuous landscapes. Satellite instruments observe radiance fields over gridded domains. Seismic arrays measure subsurface properties at scattered boreholes. In each case, neighboring observations share more information than distant ones, a property called \emph{spatial dependence} or \emph{autocorrelation}, which violates the independence assumption underlying classical statistical methods and must be explicitly modeled to produce valid inference and efficient prediction.

Gaussian process (GP) models have emerged as the dominant framework for principled spatial prediction precisely because they encode dependence structure in a flexible, mathematically coherent way. A GP defines a probability distribution over functions such that any finite collection of function values follows a multivariate Gaussian distribution \cite{rasmussen2006gaussian}. The dependence between observations is captured by a \emph{covariance function} (or \emph{kernel}), whose parameters encode spatial range, variance, and smoothness. Given observed data, the GP posterior yields analytical expressions for predictive means and variances at unobserved locations, naturally propagating uncertainty throughout the spatial domain.

In the geostatistics literature, GP-based spatial prediction is called \emph{kriging}, named after the South African mining engineer Danie Krige \cite{krige1951statistical}. Kriging and its variants have been foundational tools in geology, hydrology, and atmospheric science for decades \cite{cressie1993statistics}. More recently, the machine learning community rediscovered GPs as nonparametric Bayesian regressors, leading to significant advances in covariance function design, approximate inference, and scalable computation \cite{rasmussen2006gaussian,quinonero2005unifying}.

This survey is organized around four interconnected themes central to modern data science with dependent data. First, we examine how spatial dependencies manifest in structured datasets and how they can be quantified through variograms and correlation functions. Second, we discuss GP modeling as the appropriate framework for dependent spatial data, covering covariance function design and model selection. Third, we review prediction algorithms---kriging equations and their generalizations---that account for spatial structure to produce optimal interpolants. Fourth, we address the recovery of latent spatial structure from complex, possibly noisy or indirect observations through hierarchical and approximate GP methods.

The remainder of this paper is organized as follows. Section~\ref{sec:prelim} establishes mathematical notation and foundational concepts. Section~\ref{sec:theory} develops the core theory of GPs and spatial covariance. Section~\ref{sec:algorithms} surveys prediction algorithms and computational methods. Section~\ref{sec:examples} presents worked numerical examples. Section~\ref{sec:applications} reviews real-world applications. Section~\ref{sec:limitations} discusses limitations and assumptions. Section~\ref{sec:related} connects GPs to related methods. Section~\ref{sec:conclusion} concludes.

%=============================================================================
\section{Mathematical Preliminaries}
\label{sec:prelim}
%=============================================================================

\subsection{Notation}

We denote scalar quantities in italic ($x$, $y$, $\sigma^2$), vectors in bold lowercase ($\mathbf{x}$, $\mathbf{f}$), and matrices in bold uppercase ($\mathbf{K}$, $\mathbf{X}$). The spatial domain is $\mathcal{D} \subseteq \mathbb{R}^d$, typically $d = 2$ for planar geography. A spatial location is $\mathbf{s} \in \mathcal{D}$. Observations are denoted $\{y(\mathbf{s}_i)\}_{i=1}^n$ at locations $\{\mathbf{s}_i\}_{i=1}^n$. The notation $\mathbf{y} = (y(\mathbf{s}_1), \ldots, y(\mathbf{s}_n))^\top$ collects observations into a column vector. We write $\mathcal{N}(\bm{\mu}, \bm{\Sigma})$ for a multivariate normal distribution with mean $\bm{\mu}$ and covariance matrix $\bm{\Sigma}$.

\subsection{Random Fields}

\begin{definition}[Random Field]
A \emph{random field} (or \emph{stochastic process}) on $\mathcal{D}$ is a collection of random variables $\{Z(\mathbf{s}) : \mathbf{s} \in \mathcal{D}\}$ defined on a common probability space $(\Omega, \mathcal{F}, P)$.
\end{definition}

The \emph{mean function} $\mu(\mathbf{s}) = \mathbb{E}[Z(\mathbf{s})]$ and \emph{covariance function} $C(\mathbf{s}, \mathbf{s}') = \text{Cov}(Z(\mathbf{s}), Z(\mathbf{s}'))$ fully characterize a Gaussian random field. A covariance function must be \emph{positive semidefinite}: for any finite set of locations $\{\mathbf{s}_i\}$ and real weights $\{a_i\}$,
\begin{equation}
  \sum_{i,j} a_i a_j C(\mathbf{s}_i, \mathbf{s}_j) \geq 0.
  \label{eq:psd}
\end{equation}

\subsection{Stationarity}

\begin{definition}[Strict Stationarity]
A random field $Z$ is \emph{strictly stationary} if for any finite collection of locations $\{\mathbf{s}_i\}$ and any translation $\mathbf{h}$, the joint distribution of $\{Z(\mathbf{s}_i + \mathbf{h})\}$ equals that of $\{Z(\mathbf{s}_i)\}$.
\end{definition}

\begin{definition}[Second-Order (Weak) Stationarity]
$Z$ is \emph{second-order stationary} if $\mu(\mathbf{s})$ is constant and $C(\mathbf{s}, \mathbf{s}') = C(\mathbf{s} - \mathbf{s}')$ depends only on the lag $\mathbf{h} = \mathbf{s} - \mathbf{s}'$.
\end{definition}

\begin{definition}[Isotropy]
A stationary field is \emph{isotropic} if $C(\mathbf{h}) = C(\|\mathbf{h}\|)$ depends only on the Euclidean distance $r = \|\mathbf{h}\|$.
\end{definition}

Isotropy is a common simplifying assumption in spatial modeling, though it can be relaxed to \emph{geometric anisotropy} by applying a linear transformation to the lag vector.

\subsection{The Variogram}

The \emph{variogram} $2\gamma(\mathbf{h}) = \text{Var}(Z(\mathbf{s}+\mathbf{h}) - Z(\mathbf{s}))$ characterizes spatial dependence without requiring a finite variance, making it useful for intrinsically stationary processes.

\begin{definition}[Intrinsic Stationarity]
$Z$ is \emph{intrinsically stationary} if $\mathbb{E}[Z(\mathbf{s}+\mathbf{h}) - Z(\mathbf{s})] = 0$ and $\text{Var}(Z(\mathbf{s}+\mathbf{h}) - Z(\mathbf{s})) = 2\gamma(\mathbf{h})$ for all $\mathbf{s}, \mathbf{h}$.
\end{definition}

For second-order stationary processes, $\gamma(\mathbf{h}) = C(\mathbf{0}) - C(\mathbf{h})$. Key variogram parameters are the \emph{sill} (asymptotic variance $\sigma^2$), the \emph{range} (distance at which spatial correlation becomes negligible), and the \emph{nugget} (microscale variance or measurement error at distance zero).

%=============================================================================
\section{Core Theory}
\label{sec:theory}
%=============================================================================

\subsection{Gaussian Process Definition}

\begin{definition}[Gaussian Process]
A \emph{Gaussian process} is a collection of random variables, any finite subset of which has a joint Gaussian distribution. A GP is completely specified by its mean function $m(\mathbf{s}) = \mathbb{E}[f(\mathbf{s})]$ and covariance function $k(\mathbf{s}, \mathbf{s}') = \text{Cov}(f(\mathbf{s}), f(\mathbf{s}'))$, written:
\begin{equation}
  f(\cdot) \sim \mathcal{GP}(m(\cdot), k(\cdot, \cdot)).
  \label{eq:gp_def}
\end{equation}
\end{definition}

The GP defines a prior over functions. Given training data $\mathbf{y}$ observed at $\mathbf{S} = \{\mathbf{s}_i\}_{i=1}^n$, the posterior GP at test locations $\mathbf{S}_* = \{\mathbf{s}_j^*\}_{j=1}^m$ follows from Gaussian conditioning. Assuming an observation model $y(\mathbf{s}_i) = f(\mathbf{s}_i) + \varepsilon_i$ with $\varepsilon_i \sim \mathcal{N}(0, \sigma_n^2)$, the noisy observations satisfy:
\begin{equation}
  \mathbf{y} \sim \mathcal{N}(\mathbf{m}, \mathbf{K}_{nn} + \sigma_n^2 \mathbf{I}),
  \label{eq:marginal_lik}
\end{equation}
where $[\mathbf{K}_{nn}]_{ij} = k(\mathbf{s}_i, \mathbf{s}_j)$ and $\mathbf{m}_i = m(\mathbf{s}_i)$.

\subsection{GP Posterior and Predictive Distribution}

The joint prior over $(\mathbf{f}, \mathbf{f}_*)$ where $\mathbf{f} = f(\mathbf{S})$ and $\mathbf{f}_* = f(\mathbf{S}_*)$ is:
\begin{equation}
\begin{pmatrix} \mathbf{y} \\ \mathbf{f}_* \end{pmatrix}
\sim \mathcal{N}\!\left(
\begin{pmatrix} \mathbf{m} \\ \mathbf{m}_* \end{pmatrix},
\begin{pmatrix} \mathbf{K}_{nn} + \sigma_n^2\mathbf{I} & \mathbf{K}_{n*} \\ \mathbf{K}_{*n} & \mathbf{K}_{**} \end{pmatrix}
\right).
\label{eq:joint_prior}
\end{equation}
Conditioning on $\mathbf{y}$ yields the predictive posterior:
\begin{align}
  \bar{\mathbf{f}}_* &= \mathbf{m}_* + \mathbf{K}_{*n}(\mathbf{K}_{nn} + \sigma_n^2\mathbf{I})^{-1}(\mathbf{y} - \mathbf{m}), \label{eq:pred_mean}\\
  \mathbb{V}[\mathbf{f}_*] &= \mathbf{K}_{**} - \mathbf{K}_{*n}(\mathbf{K}_{nn} + \sigma_n^2\mathbf{I})^{-1}\mathbf{K}_{n*}. \label{eq:pred_var}
\end{align}
These equations constitute the \emph{kriging} predictor in spatial statistics. Equation~\eqref{eq:pred_mean} is the best linear unbiased predictor (BLUP) under the GP model, and \eqref{eq:pred_var} quantifies prediction uncertainty.

\subsection{Covariance Functions for Spatial Data}

The covariance function encodes the spatial dependence structure. Valid (positive definite) covariance functions for isotropic fields include the following widely used families, summarized in Table~\ref{tab:kernels}.

\begin{table}[t]
\caption{Common Isotropic Covariance Functions for Spatial GPs}
\label{tab:kernels}
\centering
\renewcommand{\arraystretch}{1.3}
\begin{tabular}{@{}llp{3.8cm}@{}}
\toprule
\textbf{Name} & \textbf{Formula} $k(r)$ & \textbf{Properties} \\
\midrule
Squared Exp. & $\sigma^2 \exp\!\left(-\tfrac{r^2}{2\ell^2}\right)$ & Infinitely differentiable; too smooth for most spatial data \\
\addlinespace
Mat\'{e}rn $\nu=\tfrac{1}{2}$ & $\sigma^2 \exp\!\left(-\tfrac{r}{\ell}\right)$ & Exponential; continuous but not differentiable \\
\addlinespace
Mat\'{e}rn $\nu=\tfrac{3}{2}$ & $\sigma^2\!\left(1+\tfrac{\sqrt{3}r}{\ell}\right)\exp\!\left(-\tfrac{\sqrt{3}r}{\ell}\right)$ & Once differentiable \\
\addlinespace
Mat\'{e}rn $\nu=\tfrac{5}{2}$ & $\sigma^2\!\left(1+\tfrac{\sqrt{5}r}{\ell}+\tfrac{5r^2}{3\ell^2}\right)\exp\!\left(-\tfrac{\sqrt{5}r}{\ell}\right)$ & Twice differentiable; widely recommended \\
\addlinespace
Rational Quad. & $\sigma^2\!\left(1+\tfrac{r^2}{2\alpha\ell^2}\right)^{-\alpha}$ & Scale mixture of squared exponential kernels \\
\addlinespace
Power Exp. & $\sigma^2 \exp\!\left(-\left(\tfrac{r}{\ell}\right)^\kappa\right)$, $\kappa\in(0,2]$ & Flexible smoothness; $\kappa=1$ exponential, $\kappa=2$ squared exp. \\
\bottomrule
\end{tabular}
\end{table}

\textbf{The Mat\'{e}rn Family.} The Mat\'{e}rn covariance \cite{matern1960spatial} is parameterized by a smoothness parameter $\nu > 0$:
\begin{equation}
  k_\nu(r) = \sigma^2 \frac{2^{1-\nu}}{\Gamma(\nu)}\left(\frac{\sqrt{2\nu}\,r}{\ell}\right)^\nu K_\nu\!\left(\frac{\sqrt{2\nu}\,r}{\ell}\right),
  \label{eq:matern}
\end{equation}
where $K_\nu$ is the modified Bessel function of the second kind and $\ell > 0$ is the length-scale. A GP with Mat\'{e}rn-$\nu$ kernel has sample paths that are $\lceil\nu\rceil - 1$ times mean-square differentiable. The Mat\'{e}rn family is the standard recommendation for spatial applications \cite{stein1999interpolation} because the squared exponential kernel's infinite differentiability is rarely justified by physical processes.

\textbf{Spectral Representation.} By Bochner's theorem \cite{bochner1933monotone}, any continuous, bounded, stationary covariance function on $\mathbb{R}^d$ has the representation:
\begin{equation}
  k(\mathbf{h}) = \int_{\mathbb{R}^d} e^{i\bm{\omega}^\top \mathbf{h}} S(\bm{\omega})\, d\bm{\omega},
  \label{eq:bochner}
\end{equation}
where $S(\bm{\omega}) \geq 0$ is the \emph{spectral density} (power spectrum). This representation underlies the \emph{spectral mixture} kernels of Wilson and Adams \cite{wilson2013gaussian}, which construct flexible kernels by modeling $S$ as a mixture of Gaussians.

\subsection{Hyperparameter Estimation via Marginal Likelihood}

GP hyperparameters $\bm{\theta} = (\sigma^2, \ell, \sigma_n^2, \ldots)$ are estimated by maximizing the log marginal likelihood (type-II maximum likelihood):
\begin{align}
  \log p(\mathbf{y} | \mathbf{X}, \bm{\theta}) &= -\frac{1}{2}\mathbf{y}^\top \mathbf{C}_\theta^{-1} \mathbf{y} - \frac{1}{2}\log|\mathbf{C}_\theta| - \frac{n}{2}\log 2\pi,
  \label{eq:log_mlik}
\end{align}
where $\mathbf{C}_\theta = \mathbf{K}_{nn}(\bm{\theta}) + \sigma_n^2 \mathbf{I}$. The gradient with respect to $\bm{\theta}_j$ is:
\begin{equation}
  \frac{\partial}{\partial \theta_j}\log p(\mathbf{y}|\mathbf{X},\bm{\theta}) = \frac{1}{2}\text{tr}\!\left(\left(\bm{\alpha}\bm{\alpha}^\top - \mathbf{C}_\theta^{-1}\right)\frac{\partial \mathbf{C}_\theta}{\partial \theta_j}\right),
  \label{eq:mlik_grad}
\end{equation}
where $\bm{\alpha} = \mathbf{C}_\theta^{-1}\mathbf{y}$. This gradient is used with gradient-based optimizers such as L-BFGS. The log marginal likelihood automatically penalizes model complexity, implementing Occam's razor without cross-validation \cite{rasmussen2006gaussian}.

\subsection{Kriging as Best Linear Unbiased Prediction}

\begin{theorem}[Simple Kriging]
Under a Gaussian process model with known mean $m(\cdot)$ and covariance $k(\cdot,\cdot)$, the minimum mean squared error linear predictor of $f(\mathbf{s}_*)$ given $\mathbf{y}$ is:
\begin{equation}
  \hat{f}(\mathbf{s}_*) = m(\mathbf{s}_*) + \mathbf{k}_*^\top \mathbf{C}^{-1}(\mathbf{y} - \mathbf{m}),
  \label{eq:simple_kriging}
\end{equation}
where $\mathbf{k}_* = (k(\mathbf{s}_*, \mathbf{s}_1), \ldots, k(\mathbf{s}_*, \mathbf{s}_n))^\top$, $\mathbf{m} = (m(\mathbf{s}_1),\ldots,m(\mathbf{s}_n))^\top$, and $\mathbf{C} = \mathbf{K}_{nn} + \sigma_n^2\mathbf{I}$. The kriging variance is:
\begin{equation}
  \sigma^2_k(\mathbf{s}_*) = k(\mathbf{s}_*, \mathbf{s}_*) - \mathbf{k}_*^\top \mathbf{C}^{-1} \mathbf{k}_*.
  \label{eq:kriging_var}
\end{equation}
\end{theorem}

\begin{proof}
This follows directly from the standard formula for conditional multivariate Gaussians; see \cite{cressie1993statistics,rasmussen2006gaussian}.
\end{proof}

\textbf{Ordinary and Universal Kriging.} When the mean is unknown but assumed constant ($m(\mathbf{s}) = \mu$), \emph{ordinary kriging} estimates $\mu$ from the data via an unbiasedness constraint, yielding weights $\bm{\lambda} = \mathbf{C}^{-1}(\mathbf{k}_* + \mathbf{1}\lambda_0)$ where $\lambda_0$ is a Lagrange multiplier enforcing $\mathbf{1}^\top\bm{\lambda} = 1$. \emph{Universal kriging} extends this to a general mean function $m(\mathbf{s}) = \mathbf{f}(\mathbf{s})^\top\bm{\beta}$, where $\mathbf{f}(\mathbf{s})$ is a vector of trend basis functions \cite{cressie1993statistics}.

%=============================================================================
\section{Algorithms and Methods}
\label{sec:algorithms}
%=============================================================================

\subsection{Standard GP Prediction Algorithm}

The standard GP prediction procedure is $O(n^3)$ in time and $O(n^2)$ in memory due to Cholesky factorization of $\mathbf{C}$.

\begin{description}
\item[Algorithm 1: GP Regression with Noise]~
\begin{algorithmic}[1]
\STATE \textbf{Input:} Training locations $\mathbf{S}$, observations $\mathbf{y}$, test locations $\mathbf{S}_*$, kernel $k(\cdot,\cdot)$, noise $\sigma_n^2$.
\STATE Compute covariance matrix $\mathbf{K}_{nn}$ with $[\mathbf{K}_{nn}]_{ij} = k(\mathbf{s}_i, \mathbf{s}_j)$.
\STATE Form $\mathbf{C} = \mathbf{K}_{nn} + \sigma_n^2\mathbf{I}$.
\STATE Compute Cholesky factorization $\mathbf{C} = \mathbf{L}\mathbf{L}^\top$.
\STATE Solve $\bm{\alpha} = \mathbf{L}^\top \backslash (\mathbf{L} \backslash (\mathbf{y} - \mathbf{m}))$.
\STATE Compute cross-covariance $\mathbf{K}_{*n}$ with $[\mathbf{K}_{*n}]_{ij} = k(\mathbf{s}_i^*, \mathbf{s}_j)$.
\STATE Predictive mean: $\bar{\mathbf{f}}_* = \mathbf{m}_* + \mathbf{K}_{*n}\bm{\alpha}$.
\STATE Solve $\mathbf{V} = \mathbf{L} \backslash \mathbf{K}_{n*}$.
\STATE Predictive variance: $\text{diag}(\mathbb{V}[\mathbf{f}_*]) = \text{diag}(\mathbf{K}_{**}) - \text{diag}(\mathbf{V}^\top \mathbf{V})$.
\STATE Log marginal likelihood: $\log p = -\frac{1}{2}\mathbf{y}^\top\bm{\alpha} - \sum_i \log L_{ii} - \frac{n}{2}\log 2\pi$.
\STATE \textbf{Output:} $\bar{\mathbf{f}}_*$, predictive variances, log marginal likelihood.
\end{algorithmic}
\end{description}

\subsection{Sparse GP Approximations}

For large $n$, exact GP inference is prohibitive. Two families of approximation dominate the literature.

\textbf{Inducing-Point Methods.} The Fully Independent Training Conditional (FITC) approximation \cite{snelson2006sparse} and the Variational Free Energy (VFE) method \cite{titsias2009variational} introduce $m \ll n$ \emph{inducing points} $\mathbf{Z} = \{\mathbf{z}_j\}_{j=1}^m$ to summarize the data. The variational posterior on $\mathbf{f}$ given inducing variables $\mathbf{u} = f(\mathbf{Z})$ is:
\begin{equation}
  q(\mathbf{f}) = \int p(\mathbf{f}|\mathbf{u}) q(\mathbf{u})\, d\mathbf{u},
  \label{eq:var_posterior}
\end{equation}
where $q(\mathbf{u}) = \mathcal{N}(\mathbf{m}_u, \mathbf{S})$ is optimized by maximizing the evidence lower bound (ELBO):
\begin{align}
  \mathcal{L} &= \sum_{i=1}^n \mathbb{E}_{q(f_i)}\!\left[\log p(y_i | f_i)\right] - \mathrm{KL}[q(\mathbf{u}) \| p(\mathbf{u})].
  \label{eq:elbo}
\end{align}
This reduces computation to $O(nm^2)$ per iteration. The inducing locations can be optimized jointly with hyperparameters.

\begin{description}
\item[Algorithm 2: Sparse Variational GP (SVGP) Training]~
\begin{algorithmic}[1]
\STATE \textbf{Input:} Data $\{(\mathbf{s}_i, y_i)\}$, $m$ inducing locations $\mathbf{Z}$ (initialized by k-means), kernel $k$, likelihood $p(y|f)$.
\STATE Initialize variational parameters $\mathbf{m}_u, \mathbf{S}$ (e.g., $\mathbf{m}_u = \mathbf{0}$, $\mathbf{S} = \mathbf{K}_{mm}$).
\REPEAT
  \STATE Sample mini-batch $\mathcal{B} \subset \{1,\ldots,n\}$.
  \STATE Compute $\mathbf{K}_{mm}$, $\mathbf{K}_{bm}$ for batch indices $\mathcal{B}$.
  \STATE Compute $q(f_i)$ for $i \in \mathcal{B}$ via $q(f_i) = \mathcal{N}(\mu_i, v_i)$:
      \begin{align*}
        \mu_i &= \mathbf{k}_{im}\mathbf{K}_{mm}^{-1}\mathbf{m}_u, \\
        v_i &= k_{ii} - \mathbf{k}_{im}\mathbf{K}_{mm}^{-1}(\mathbf{K}_{mm}-\mathbf{S})\mathbf{K}_{mm}^{-1}\mathbf{k}_{mi}.
      \end{align*}
  \STATE Estimate ELBO (scale batch term by $n/|\mathcal{B}|$).
  \STATE Update $\mathbf{m}_u$, $\mathbf{S}$, $\mathbf{Z}$, $\bm{\theta}$ via Adam gradient step.
\UNTIL{convergence}
\STATE \textbf{Output:} Optimized $\mathbf{m}_u$, $\mathbf{S}$, $\mathbf{Z}$, $\bm{\theta}$.
\end{algorithmic}
\end{description}

\textbf{Gaussian Markov Random Field (GMRF) Approximations.} Lindgren et al.\ \cite{lindgren2011explicit} showed that solutions to the stochastic partial differential equation (SPDE)
\begin{equation}
  (\kappa^2 - \Delta)^{\alpha/2}(\tau Z(\mathbf{s})) = \mathcal{W}(\mathbf{s}),
  \label{eq:spde}
\end{equation}
where $\mathcal{W}$ is a Gaussian white noise process and $\Delta$ is the Laplacian, have Mat\'{e}rn covariance with $\nu = \alpha - d/2$. Finite element discretization of \eqref{eq:spde} yields a sparse precision matrix, enabling $O(n^{3/2})$ computation in $d=2$ via sparse Cholesky methods \cite{rue2005gaussian}.

\subsection{Kriging Variants}

Table~\ref{tab:kriging_variants} summarizes the main kriging variants and their assumptions.

\begin{table}[t]
\caption{Kriging Variants and Assumptions}
\label{tab:kriging_variants}
\centering
\renewcommand{\arraystretch}{1.3}
\begin{tabular}{@{}lll@{}}
\toprule
\textbf{Method} & \textbf{Mean Assumption} & \textbf{Key Feature} \\
\midrule
Simple Kriging & Known constant $\mu$ & Textbook BLUP \\
Ordinary Kriging & Unknown constant & Unbiasedness constraint \\
Universal Kriging & $\mathbf{f}(\mathbf{s})^\top\bm{\beta}$, $\bm{\beta}$ unknown & Linear trend removal \\
Co-Kriging & Multivariate $Z$ & Uses auxiliary variables \\
Indicator Kriging & Threshold indicators & Probabilistic maps \\
Trans-Gaussian & Transformed data & Lognormal, Box-Cox \\
Disjunctive & Nonlinear $f(Z)$ & Non-Gaussian goals \\
Bayesian Kriging & Full Bayesian & Uncertainty in $\bm{\theta}$ \\
\bottomrule
\end{tabular}
\end{table}

\subsection{Variogram Estimation and Fitting}

Before kriging, the variogram must be estimated from data and a valid parametric model fit. The \emph{method of moments} (empirical variogram) estimator is:
\begin{equation}
  \hat{\gamma}(h) = \frac{1}{2|N(h)|}\sum_{(i,j)\in N(h)} (y(\mathbf{s}_i) - y(\mathbf{s}_j))^2,
  \label{eq:empirical_variogram}
\end{equation}
where $N(h) = \{(i,j) : \|\mathbf{s}_i - \mathbf{s}_j\| \approx h\}$ collects pairs at distance $h$ (within a tolerance band). A parametric model $\gamma(h;\bm{\theta})$ is then fit by weighted least squares or maximum likelihood to ensure validity.

\textbf{Model selection} between competing covariance families can use the log marginal likelihood \eqref{eq:log_mlik}, cross-validation, or information criteria such as AIC and BIC. Stein \cite{stein1999interpolation} showed that the behavior of the covariance near the origin (smoothness) is most critical for interpolation, motivating careful choice of $\nu$ in the Mat\'{e}rn family.

%=============================================================================
\section{Worked Examples}
\label{sec:examples}
%=============================================================================

\subsection{Example 1: 1-D Spatial Interpolation}

Consider $n = 5$ noiseless observations at locations $\mathbf{s} = (0, 1, 2, 4, 5)^\top$ with values $\mathbf{y} = (0.2, 0.8, 0.6, -0.3, -0.5)^\top$. We fit a GP with Mat\'{e}rn-$\tfrac{3}{2}$ kernel, $\sigma^2 = 1$, $\ell = 1.5$, and no noise ($\sigma_n^2 = 0$). The covariance matrix $\mathbf{K}$ is computed via \eqref{eq:matern}. For example, the $(1,2)$ entry with $r = 1$, $\ell = 1.5$ is:
\begin{equation}
  k(1) = 1\cdot\!\left(1 + \frac{\sqrt{3}}{1.5}\right)\exp\!\left(-\frac{\sqrt{3}}{1.5}\right) \approx 0.569.
\end{equation}
Solving the kriging system \eqref{eq:simple_kriging}--\eqref{eq:kriging_var} at a test point $s_* = 3$ gives predictive mean $\hat{f}(3) \approx 0.14$ and standard deviation $\sigma_k(3) \approx 0.41$, reflecting the moderate distance from observations and uncertainty about the function value in the data gap.

The kriging predictor interpolates exactly at observed locations when $\sigma_n^2 = 0$: $\hat{f}(\mathbf{s}_i) = y(\mathbf{s}_i)$ for all $i$. The predictive standard deviation is zero at observed locations and increases away from them, reaching its maximum in data-sparse regions. This illustrates the spatial uncertainty structure that distinguishes GP prediction from classical regression.

\subsection{Example 2: Variogram Fitting for Temperature Data}

Suppose we have daily maximum temperatures (in $^\circ$C) at $n = 150$ weather stations across a $500 \times 500$ km region. After removing a linear trend $m(\mathbf{s}) = \beta_0 + \beta_1 s_1 + \beta_2 s_2$ fit by ordinary least squares, we compute the empirical variogram of the residuals at lag classes $h \in \{25, 50, 75, \ldots, 250\}$ km using \eqref{eq:empirical_variogram}.

The empirical variogram rises from a nugget of approximately $\hat{\gamma}(0^+) \approx 0.8$ (interpreted as measurement error plus microscale variability) to a sill of $\hat{\gamma}(\infty) \approx 6.2$ at a practical range of approximately $180$ km. We fit a Mat\'{e}rn-$\tfrac{5}{2}$ model by weighted least squares minimizing:
\begin{equation}
  \sum_h \frac{|N(h)|}{2}\left(\frac{\hat{\gamma}(h)}{\gamma(h;\bm{\theta})} - 1\right)^2,
  \label{eq:wls_vario}
\end{equation}
obtaining estimates $\hat{\sigma}^2 = 5.4$, $\hat{\ell} = 78$ km, and $\hat{\sigma}_n^2 = 0.8$. Universal kriging with this fitted model is then used to produce a continuous temperature map over the spatial domain, with the kriging standard deviation map revealing higher uncertainty in regions with sparse station coverage.

\subsection{Example 3: Latent GP for Spatial Count Data}

Let $y(\mathbf{s}_i) \sim \text{Poisson}(\exp(f(\mathbf{s}_i)))$ represent disease counts in $n = 200$ census tracts, where $f \sim \mathcal{GP}(0, k)$ is a latent spatial log-intensity. The non-Gaussian likelihood breaks the analytical tractability of standard kriging. We apply Laplace approximation to the posterior $p(\mathbf{f}|\mathbf{y}) \propto p(\mathbf{y}|\mathbf{f})p(\mathbf{f})$ by finding the mode $\hat{\mathbf{f}}$ via Newton's method:
\begin{equation}
  \hat{\mathbf{f}} \leftarrow \hat{\mathbf{f}} + (\mathbf{K}^{-1} + \mathbf{W})^{-1}(\nabla \log p(\mathbf{y}|\hat{\mathbf{f}}) - \mathbf{K}^{-1}\hat{\mathbf{f}}),
  \label{eq:newton_laplace}
\end{equation}
where $\mathbf{W} = -\nabla^2 \log p(\mathbf{y}|\hat{\mathbf{f}})$ is the diagonal Hessian. For Poisson, $\mathbf{W} = \text{diag}(\exp(\hat{f}_i))$. The Laplace approximation $q(\mathbf{f}) = \mathcal{N}(\hat{\mathbf{f}}, (\mathbf{K}^{-1} + \mathbf{W})^{-1})$ then propagates spatial dependence to the predicted count surface, revealing geographic clustering of disease risk that would be obscured by model-free smoothing.

%=============================================================================
\section{Applications in Data Science}
\label{sec:applications}
%=============================================================================

\subsection{Environmental Monitoring and Interpolation}

One of the canonical applications of spatial GPs is the interpolation of environmental variables from sparse sensor networks. The US Environmental Protection Agency and European Environment Agency routinely use kriging to produce national air quality maps from regulatory monitor networks \cite{hoek2008review}. Spatial GPs are used to estimate PM$_{2.5}$, ozone, and nitrogen dioxide concentrations at unmonitored locations, combining ground monitor data with remote sensing covariates through co-kriging or regression-kriging \cite{knotters1995comparison}. The resulting surfaces inform regulatory exposure assessments and epidemiological studies.

GP methods have also been central to spatial climate downscaling, where coarse-resolution climate model output is adjusted to the scale of local observations \cite{willmott1985statistics}. By modeling the residual between the coarse model and local station data as a GP, practitioners recover fine-scale spatial variability that climate models cannot resolve.

\subsection{Geostatistics and Natural Resource Estimation}

Kriging originated in mining geostatistics and remains widely used for ore grade estimation, groundwater modeling, and petroleum reservoir characterization. In these applications, the GP provides not only a point estimate of the resource concentration but a full probability distribution, enabling risk-weighted decision-making about extraction \cite{cressie1993statistics}. Object-based geostatistical simulation using conditional GP realizations produces multiple equiprobable spatial scenarios, quantifying uncertainty in reserve estimates.

\subsection{Bayesian Optimization}

Bayesian optimization (BO) uses a GP surrogate model to sequentially optimize expensive black-box functions \cite{snoek2012practical}. The acquisition function---expected improvement, upper confidence bound, or probability of improvement---uses the GP predictive mean and variance to balance exploration and exploitation. In hyperparameter tuning for machine learning pipelines, BO with GP surrogates consistently outperforms grid and random search, particularly in high-dimensional settings where function evaluations are costly. Spatial structure arises when the black-box function exhibits locality in hyperparameter space, which GP kernels naturally encode.

\subsection{Gaussian Process Regression in Machine Learning}

In the broader machine learning context, GP regression is used for nonparametric function estimation with uncertainty quantification. Applications include time series forecasting \cite{roberts2013gaussian}, recommender systems, and scientific emulation of expensive simulators \cite{santner2003design}. The ability to provide calibrated predictive intervals distinguishes GPs from neural network regressors in safety-critical settings such as autonomous vehicles and medical diagnosis.

\subsection{Spatial Epidemiology and Disease Mapping}

Latent GP models for areal and point-referenced disease data identify spatial clusters of elevated risk beyond what chance variation would predict \cite{besag1991bayesian}. The Besag-York-Molli\'{e} (BYM) model combines an intrinsic conditional autoregressive (ICAR) prior---a discrete GP analogue---with an unstructured random effect to separate spatially structured from unstructured overdispersion. These models are estimated within the integrated nested Laplace approximation (INLA) framework \cite{rue2009approximate}, enabling scalable Bayesian inference for large spatial datasets.

\subsection{Remote Sensing and Image Analysis}

GP models applied to satellite imagery exploit both spatial and spectral dependence. Multi-output GPs model correlations between spectral bands through shared latent processes \cite{alvarez2012kernels}, improving retrieval accuracy for biophysical parameters such as leaf area index and soil moisture. Spatial GPs also underlie kriging-based gap-filling for cloud-contaminated satellite imagery, where missing pixels are predicted from neighboring clear-sky observations.

%=============================================================================
\section{Limitations and Assumptions}
\label{sec:limitations}
%=============================================================================

\subsection{Stationarity Assumption}

The standard GP model assumes stationarity: the covariance between two points depends only on their separation vector, not their absolute locations. This assumption is violated when:
\begin{itemize}
  \item The spatial process has heterogeneous smoothness (e.g., terrain that is smooth in valleys but rough over mountain ridges).
  \item The range or variance differs systematically across subregions.
  \item Physical barriers (coastlines, mountain ranges) break spatial continuity.
\end{itemize}
Nonstationary extensions include locally stationary models \cite{fuentes2001new}, deformation approaches that map the spatial domain to an isotropic reference space \cite{sampson1992nonparametric}, and deep kernel learning where the feature map $\mathbf{s} \mapsto \phi(\mathbf{s})$ is learned by a deep neural network \cite{wilson2016deep}.

\subsection{Computational Complexity}

Standard GP inference scales as $O(n^3)$ in time and $O(n^2)$ in storage. For $n \geq 10{,}000$ locations, exact computation becomes infeasible on typical hardware. Although sparse approximations (Section~\ref{sec:algorithms}) reduce complexity, they introduce approximation error and require careful selection of inducing point count and placement. The SPDE approach achieves near-linear scaling but is tied to Mat\'{e}rn covariances on Euclidean or mesh domains.

\subsection{Gaussian Assumption}

GP models assume Gaussian marginals, which may be inappropriate for:
\begin{itemize}
  \item Skewed or heavy-tailed data (precipitation totals, financial returns).
  \item Binary or categorical spatial data (land cover classification).
  \item Count data (disease surveillance, species abundance).
\end{itemize}
Transformation approaches (lognormal kriging, Box-Cox transforms) partially address skewness. For fundamentally non-Gaussian responses, latent GP models with non-Gaussian likelihoods (Section~\ref{sec:examples}, Example 3) require approximate inference.

\subsection{Hyperparameter Sensitivity}

GP predictions can be sensitive to the choice of covariance function and its hyperparameters, particularly when $n$ is small or the observations are noisy. Maximum likelihood estimation of hyperparameters can converge to local optima, and the marginal likelihood surface is often multimodal. Cross-validation strategies and multiple restarts from diverse initial conditions help mitigate this issue \cite{rasmussen2006gaussian}. Full Bayesian treatment of hyperparameters via MCMC is possible but computationally demanding.

\subsection{Isotropy and Low-Dimensional Inputs}

Standard GP kernels scale poorly with input dimension $d$ due to the curse of dimensionality: all points become equidistant in high dimensions, causing the correlation structure to degenerate. GP spatial models are most natural and well-calibrated when $d \leq 3$. Additive GPs, cylindrical kernels for directional data, and deep GP architectures extend the framework to higher dimensions.

%=============================================================================
\section{Connections to Related Methods}
\label{sec:related}
%=============================================================================

\subsection{Gaussian Markov Random Fields}

GMRFs \cite{rue2005gaussian} define spatial dependence through a sparse precision (inverse covariance) matrix $\mathbf{Q}$, in contrast to GPs which directly parameterize the covariance. GMRFs are natural for areal data (zip codes, census tracts) where the neighborhood graph defines adjacency. The SPDE connection \cite{lindgren2011explicit} shows that Mat\'{e}rn GPs are the continuous-domain limits of certain GMRFs, enabling GMRF computational methods for GP problems.

\subsection{Kernel Methods and Support Vector Regression}

GP regression with a squared exponential kernel is closely related to kernel ridge regression (KRR) with the same kernel: both yield the same predictive mean \cite{rasmussen2006gaussian}. The difference is that GPs additionally provide calibrated predictive uncertainty through the posterior variance, which KRR does not. The representer theorem guarantees that both solutions lie in the reproducing kernel Hilbert space (RKHS) associated with $k$.

\subsection{Deep Gaussian Processes}

Deep GPs \cite{damianou2013deep} compose multiple GP layers, with the output of each layer serving as the input to the next:
\begin{equation}
  f^{(L)} = \mathcal{GP}(0, k(f^{(L-1)}, f^{(L-1)})), \quad f^{(l)} = \mathcal{GP}(0, k(f^{(l-1)}, f^{(l-1)})).
  \label{eq:deep_gp}
\end{equation}
This composition induces non-Gaussian, non-stationary priors on the function, with greater expressiveness than single-layer GPs. Doubly stochastic variational inference \cite{salimbeni2017doubly} enables scalable training. Deep GPs are particularly suited to spatial modeling of heterogeneous processes.

\subsection{Neural Processes}

Neural processes (NPs) \cite{garnelo2018neural} amortize GP inference by training a neural network to map context sets (observed data) to predictive distributions at query locations. Unlike GPs, NPs do not require specifying a covariance function or solving linear systems at test time, enabling rapid adaptation in meta-learning settings. However, NPs sacrifice the theoretical guarantees of GPs and may produce less calibrated uncertainty estimates.

\subsection{Spatial Autoregressions}

Spatial autoregressive (SAR) and spatial error (SEM) models from spatial econometrics parameterize dependence through a spatial weights matrix $\mathbf{W}$ (e.g., queen contiguity):
\begin{equation}
  \mathbf{y} = \rho \mathbf{W}\mathbf{y} + \mathbf{X}\bm{\beta} + \bm{\varepsilon},
  \label{eq:sar}
\end{equation}
where $\rho$ is the spatial autoregressive coefficient. While less flexible than GPs, SAR models are computationally efficient and well-integrated into econometric inference frameworks \cite{anselin1988spatial}. They can be interpreted as a discrete approximation to GP models with exponential covariance.

\subsection{Interpolation Methods Without Probabilistic Foundation}

Inverse distance weighting (IDW), spline interpolation, and natural neighbor methods produce spatial predictions without a probabilistic model. They lack the uncertainty quantification of GPs and do not provide a mechanism for selecting the smoothness of the interpolant from data. However, they remain useful for exploratory analysis and when computational resources preclude GP fitting.

%=============================================================================
\section{Conclusion}
\label{sec:conclusion}
%=============================================================================

This survey has reviewed Gaussian process models as a unified, principled framework for spatial prediction. We established the mathematical foundations of spatial stationarity, covariance functions, and variogram modeling, and derived the kriging predictor as the optimal linear estimator under the GP model. We surveyed key algorithms---exact GP regression, sparse inducing-point methods, SPDE/GMRF approximations, and variogram fitting---and illustrated their application through worked examples spanning 1-D interpolation, temperature mapping, and latent spatial risk modeling. Real-world applications in environmental monitoring, natural resource estimation, Bayesian optimization, epidemiology, and remote sensing demonstrate the broad relevance of GP spatial models to data science.

The four themes of dependent data modeling are thoroughly instantiated in the GP framework. Spatial dependencies are identified through empirical variograms and covariance estimation. The GP model appropriately accounts for these dependencies in its prior and likelihood. Kriging equations produce optimal predictions that exploit the spatial structure, with calibrated uncertainty. Latent GP formulations recover hidden spatial fields from noisy or indirect observations, revealing structure that cannot be directly measured.

Key limitations include computational scalability, the stationarity assumption, Gaussian marginal constraints, and hyperparameter sensitivity. Active research frontiers address each of these: scalable sparse and SPDE methods reduce computational cost, nonstationary and deep GP models increase expressiveness, and robust likelihood and transformation methods handle non-Gaussian data. The integration of GP spatial models with deep learning---through deep kernels, deep GPs, and neural processes---is an especially active area with substantial potential for data science applications.

GP models occupy a privileged position in spatial statistics: they are simultaneously the framework of classical geostatistics (kriging) and modern machine learning (Bayesian nonparametrics), connecting decades of spatial theory to contemporary probabilistic machine learning. Their rigorous treatment of uncertainty, their connection to physical models through SPDEs, and their flexibility through kernel design make them indispensable tools for any data scientist working with spatial, temporal, or otherwise structured data.

%=============================================================================
% REFERENCES
%=============================================================================

\begin{thebibliography}{25}

\bibitem{rasmussen2006gaussian}
C.~E. Rasmussen and C.~K.~I. Williams, \emph{Gaussian Processes for Machine Learning}. Cambridge, MA: MIT Press, 2006.

\bibitem{krige1951statistical}
D.~G. Krige, ``A statistical approach to some basic mine valuation problems on the Witwatersrand,'' \emph{J. Chem. Metall. Min. Soc. South Africa}, vol.~52, no.~6, pp.~119--139, 1951.

\bibitem{cressie1993statistics}
N.~Cressie, \emph{Statistics for Spatial Data}, revised ed. New York: Wiley, 1993.

\bibitem{matern1960spatial}
B.~Mat\'{e}rn, \emph{Spatial Variation}, 2nd ed., Lecture Notes in Statistics, vol.~36. Berlin: Springer, 1986.

\bibitem{stein1999interpolation}
M.~L. Stein, \emph{Interpolation of Spatial Data: Some Theory for Kriging}. New York: Springer, 1999.

\bibitem{bochner1933monotone}
S.~Bochner, \emph{Lectures on Fourier Integrals}. Princeton, NJ: Princeton University Press, 1959.

\bibitem{lindgren2011explicit}
F.~Lindgren, H.~Rue, and J.~Lindstr\"{o}m, ``An explicit link between Gaussian fields and Gaussian Markov random fields: the stochastic partial differential equation approach,'' \emph{J. R. Stat. Soc. Ser. B Stat. Methodol.}, vol.~73, no.~4, pp.~423--498, 2011.

\bibitem{rue2005gaussian}
H.~Rue and L.~Held, \emph{Gaussian Markov Random Fields: Theory and Applications}. Boca Raton, FL: Chapman \& Hall/CRC, 2005.

\bibitem{quinonero2005unifying}
J.~Qui\~{n}onero-Candela and C.~E. Rasmussen, ``A unifying view of sparse approximate Gaussian process regression,'' \emph{J. Mach. Learn. Res.}, vol.~6, pp.~1939--1959, 2005.

\bibitem{snelson2006sparse}
E.~Snelson and Z.~Ghahramani, ``Sparse Gaussian processes using pseudo-inputs,'' in \emph{Advances in Neural Information Processing Systems}, vol.~18, 2006, pp.~1257--1264.

\bibitem{titsias2009variational}
M.~K. Titsias, ``Variational learning of inducing variables in sparse Gaussian processes,'' in \emph{Proc. 12th Int. Conf. Artif. Intell. Stat. (AISTATS)}, 2009, pp.~567--574.

\bibitem{rue2009approximate}
H.~Rue, S.~Martino, and N.~Chopin, ``Approximate Bayesian inference for latent Gaussian models by using integrated nested Laplace approximations,'' \emph{J. R. Stat. Soc. Ser. B Stat. Methodol.}, vol.~71, no.~2, pp.~319--392, 2009.

\bibitem{wilson2013gaussian}
A.~G. Wilson and R.~P. Adams, ``Gaussian process kernels for pattern discovery and extrapolation,'' in \emph{Proc. 30th Int. Conf. Mach. Learn. (ICML)}, 2013, pp.~1067--1075.

\bibitem{snoek2012practical}
J.~Snoek, H.~Larochelle, and R.~P. Adams, ``Practical Bayesian optimization of machine learning algorithms,'' in \emph{Advances in Neural Information Processing Systems}, vol.~25, 2012, pp.~2951--2959.

\bibitem{damianou2013deep}
A.~Damianou and N.~D. Lawrence, ``Deep Gaussian processes,'' in \emph{Proc. 16th Int. Conf. Artif. Intell. Stat. (AISTATS)}, 2013, pp.~207--215.

\bibitem{salimbeni2017doubly}
H.~Salimbeni and M.~Deisenroth, ``Doubly stochastic variational inference for deep Gaussian processes,'' in \emph{Advances in Neural Information Processing Systems}, vol.~30, 2017.

\bibitem{garnelo2018neural}
M.~Garnelo \emph{et al.}, ``Neural processes,'' \emph{arXiv preprint arXiv:1807.01622}, 2018.

\bibitem{besag1991bayesian}
J.~Besag, J.~York, and A.~Molli\'{e}, ``Bayesian image restoration, with two applications in spatial statistics,'' \emph{Ann. Inst. Stat. Math.}, vol.~43, no.~1, pp.~1--20, 1991.

\bibitem{wilson2016deep}
A.~G. Wilson, Z.~Hu, R.~Salakhutdinov, and E.~P. Xing, ``Deep kernel learning,'' in \emph{Proc. 19th Int. Conf. Artif. Intell. Stat. (AISTATS)}, 2016, pp.~370--378.

\bibitem{alvarez2012kernels}
M.~A. \'{A}lvarez, L.~Rosasco, and N.~D. Lawrence, ``Kernels for vector-valued functions: a review,'' \emph{Found. Trends Mach. Learn.}, vol.~4, no.~3, pp.~195--266, 2012.

\bibitem{santner2003design}
T.~J. Santner, B.~J. Williams, and W.~I. Notz, \emph{The Design and Analysis of Computer Experiments}. New York: Springer, 2003.

\bibitem{roberts2013gaussian}
S.~Roberts, M.~Osborne, M.~Ebden, S.~Reece, N.~Gibson, and S.~Aigrain, ``Gaussian processes for time-series modelling,'' \emph{Philos. Trans. R. Soc. A}, vol.~371, no.~1984, p.~20110550, 2013.

\bibitem{hoek2008review}
G.~Hoek \emph{et al.}, ``A review of land-use regression models to assess spatial variation of outdoor air pollution,'' \emph{Atmos. Environ.}, vol.~42, no.~33, pp.~7561--7578, 2008.

\bibitem{anselin1988spatial}
L.~Anselin, \emph{Spatial Econometrics: Methods and Models}. Dordrecht: Kluwer Academic Publishers, 1988.

\bibitem{fuentes2001new}
M.~Fuentes, ``A new high frequency kriging approach for nonstationary environmental processes,'' \emph{Environmetrics}, vol.~12, no.~5, pp.~469--483, 2001.

\end{thebibliography}

\end{document}
